\begin{remarks}[7.3.9]
\label{IV.2.s7.3.9-fr}
\begin{enumerate}
\item[(i)] Cada una de las propiedades (iv$'$),
$(\mathrm v'_n)$, (vi$'$), (vii$'$) implica respectivamente la
propiedad correspondiente (iv), $(\mathrm v_n)$, (vi), (vii). La
propiedad (iv) implica todas las propiedades de (i) a (vii), y la
propiedad (iv$'$) implica todas las propiedades enumeradas en
(7.3.8). Recordemos también que $(\mathrm i_0)$ equivale a (ii);
la conjunción de $(\mathrm{iii}_1)$ y $(\mathrm v_0)$ (resp. de
$(\mathrm{iii}_1)$ y $(\mathrm v'_0)$) equivale a (vi) (resp. a
(vi$'$)) (5.8.5); por último, la conjunción de $(\mathrm{iii}_2)$
y $(\mathrm v_1)$ (resp. de $(\mathrm{iii}_2)$ y $(\mathrm v'_1)$)
equivale a (vii) (resp. a (vii$'$)) (5.8.6).
\item[(ii)] En todos los ejemplos de (7.3.8), la propiedad
$\boldsymbol Q(\mathscr O_z,k)$ que sirve para definir
$\boldsymbol P$ es tal que, para toda \emph{generalización} $z'$ de
$z$ en $Z$, $\boldsymbol Q(\mathscr O_z,k)$ \emph{implica}
$\boldsymbol Q(\mathscr O_{z'},k)$: en virtud de (2.3.4), basta
verificarlo para las propiedades (i) a (vii) (e incluso (i) a (v),
según la observación (7.3.9, (i))); ahora bien, esto forma parte de
la definición para (iii) y (v) (5.7.2 y 5.8.2), se deduce de (6.3.9)
para (i) y, por último, de (\textbf 0, 16.5.10 y 17.3.2) para (ii) y (iv).
\end{enumerate}
\end{remarks}
